\section*{Chapitre \ref{ch:b1:cross-sections} --- Les coupes géologiques}

\begin{solution}{exo:b1:cross-sections:1}
À 1:25\,000, \qty{1}{cm} sur la carte valent \qty{250}{m} sur le terrain.
L'échelle verticale est de \qty{1}{cm} pour \qty{100}{m}, donc l'échelle verticale
est $250/100 = 2.5$ fois l'échelle horizontale~: $e = 2.5$.
\end{solution}

\begin{solution}{exo:b1:cross-sections:2}
$\gamma = \ang{90} - \ang{60} = \ang{30}$, et $\tan\beta = \tan\ang{40}
\cos\ang{30} = 0.8391\times0.8660 = 0.7267$~: $\beta = \ang{36.0}$, la couche
descendant vers l'extrémité est de la coupe. Dans la coupe d'\omterm{def:b1:spherical-astronomy:horizontal}{azimut}
\ang{000}, $\gamma = \ang{90}$ et $\beta = 0$~: la coupe suit la direction et la
couche apparaît horizontale.
\end{solution}

\begin{solution}{exo:b1:cross-sections:3}
$t = h\cos\delta = 120\times\cos\ang{12.16} = 120\times0.9776 =
\qty{117.3}{m}$.
\end{solution}

\begin{solution}{exo:b1:cross-sections:4}
$\tan\ang{12} = 0.2126$. Avec $e = 2$~: $\arctan 0.4251 = \ang{23.0}$~; avec
$e = 3$~: $\arctan 0.6377 = \ang{32.5}$~; avec $e = 5$~: $\arctan 1.0628 =
\ang{46.7}$. L'angle croît bien moins vite que $e$ lui-même, et n'atteint jamais
\ang{90}.
\end{solution}

\begin{solution}{exo:b1:cross-sections:5}
(a) $t = 300\sin\ang{25} = \qty{126.8}{m}$. (b) $t = 300\sin\ang{35} =
\qty{172.1}{m}$. (c) $t = 300\sin\ang{15} = \qty{77.6}{m}$. Une même largeur
d'\omterm{def:b1:geological-maps:geological-map}{affleurement} peut correspondre à des épaisseurs qui diffèrent de plus d'un
facteur deux~: il faut toujours tenir compte de la pente du terrain.
\end{solution}

\begin{solution}{exo:b1:cross-sections:6}
Avec le \omterm{def:b1:geological-maps:attitude}{sens du pendage} $\alpha$ et le \omterm{def:b1:geological-maps:attitude}{pendage} $\delta$, le théorème donne dans la
coupe nord--sud ($\theta = 0$) $\tan\ang{20} = \tan\delta\cos\alpha$,
et dans la coupe est--ouest ($\theta = \ang{90}$) $\tan\ang{30} =
\tan\delta\cos(\alpha - \ang{90}) = \tan\delta\sin\alpha$. En divisant,
$\tan\alpha = 0.5774/0.3640 = 1.586$, $\alpha = \ang{057.8}$~; en additionnant
les carrés, $\tan\delta = \sqrt{0.3640^2 + 0.5774^2} = 0.6825$, $\delta =
\ang{34.3}$. Les deux \omterm{def:b1:cross-sections:apparent-dip}{pendages apparents} sont inférieurs au \omterm{def:b1:geological-maps:attitude}{pendage} réel, comme
il se doit.
\end{solution}

\begin{solution}{exo:b1:cross-sections:7}
$100/\tan\ang{12.16} = 100/0.2155 = \qty{464.1}{m}$. Sur la carte, les
\omterm{def:b1:cross-sections:structure-contour}{isohypses structurales} du flanc est sont bien espacées d'un peu moins d'un
demi-kilomètre, équidistantes et parallèles~: le flanc est un plan.
\end{solution}

\begin{solution}{exo:b1:cross-sections:8}
B1~: $524.02 - 488.98 = \qty{35.04}{m}$~; B3~: $529.98 - 494.93 =
\qty{35.05}{m}$. Au-dessus de B2, la base de l'argilite est le toit du
calcaire, \qty{180}{m} au-dessus de sa base~: $280.16 + 180 = \qty{460.16}{m}$,
plus haut que le sol (\qty{380.64}{m}), ce qui explique qu'elle soit érodée à cet
endroit. Entre B1 et B2, le contact monte de $460.16 - 35.04 = \qty{425.12}{m}$
sur \qty{2000}{m}~: une pente de $0.2126$, un \omterm{def:b1:cross-sections:apparent-dip}{pendage apparent} de \ang{12.0}. B1
et B3, à la même distance de part et d'autre de B2, trouvent chaque contact à la
même altitude à \qty{1}{cm} près~: le \omterm{def:b1:geological-maps:fold}{pli} est symétrique par rapport à un plan
axial vertical en $x = \qty{3.5}{km}$, ses deux flancs plongeant également, et
aucune faille ne recoupe la coupe entre eux.
\end{solution}

\begin{solution}{exo:b1:cross-sections:9}
\omterm{def:b1:cross-sections:throw}{Rejet horizontal} $= 150/\tan\ang{50} = 150/1.1918 = \qty{125.9}{m}$. Dans la
coupe d'\omterm{def:b1:spherical-astronomy:horizontal}{azimut} \ang{150}, $\gamma = \ang{30}$ et $\tan\beta = 1.1918\times0.8660 =
1.0321$~: la faille apparaît avec le \omterm{def:b1:cross-sections:apparent-dip}{pendage apparent} \ang{45.9}. Les deux traces
d'une couche horizontale sont distantes de \qty{150}{m} à la verticale sur cette
trace, donc de $150/1.0321 = \qty{145.3}{m}$ le long de la coupe, plus que le
\omterm{def:b1:cross-sections:throw}{rejet horizontal}, puisque la coupe recoupe la faille obliquement.
\end{solution}

\begin{solution}{exo:b1:cross-sections:10}
$\cos\gamma = \tan\ang{20}/\tan\ang{40} = 0.3640/0.8391 = 0.4338$, donc
$\gamma = \ang{64.3}$~: les coupes d'\omterm{def:b1:spherical-astronomy:horizontal}{azimuts} $\alpha \pm \ang{64.3}$
(et les directions opposées, où la couche monte). Le \omterm{def:b1:cross-sections:apparent-dip}{pendage apparent} vaut
\ang{40.0}, \ang{39.6}, \ang{38.3} et \ang{36.0} pour $\gamma = 0$, \ang{10},
\ang{20}, \ang{30}, mais \ang{22.8}, \ang{16.0}, \ang{8.3} et $0$ pour
$\gamma = \ang{60}$, \ang{70}, \ang{80}, \ang{90}. En effet, en dérivant
$\tan\beta = \tan\delta\cos\gamma$,
$\dfrac{\mathrm d\beta}{\mathrm d\gamma} =
-\dfrac{\tan\delta\sin\gamma}{1 + \tan^2\delta\cos^2\gamma}$, qui est nul en
$\gamma = 0$ et vaut $-\tan\delta$ en $\gamma = \ang{90}$~: une coupe un peu
écartée du \omterm{def:b1:geological-maps:attitude}{sens du pendage} montre presque le \omterm{def:b1:geological-maps:attitude}{pendage} réel, alors que près de la
direction chaque degré d'\omterm{def:b1:spherical-astronomy:horizontal}{azimut} change le \omterm{def:b1:cross-sections:apparent-dip}{pendage apparent} de presque un degré.
\end{solution}

\begin{solution}{exo:b1:cross-sections:11}
Les couches plongent vers l'est, donc la normale à la stratification, suivie vers
le bas, est inclinée de \ang{12.16} par rapport à la verticale vers l'ouest, vers
\ang{260.7}. (a) Un sondage vertical fait avec elle l'angle $\psi = \ang{12.16}$~:
$176.0/\cos\ang{12.16} = \qty{180.0}{m}$, l'\omterm{def:b1:cross-sections:thickness}{épaisseur verticale}. (b) Un
sondage plongeant de \ang{60} vers \ang{260.7} est incliné de \ang{30} par
rapport à la verticale, du même côté que la normale~: $\psi = \ang{30} - \ang{12.16} =
\ang{17.84}$, $\cos\psi = 0.9519$, et la carotte mesure $176.0/0.9519 =
\qty{184.9}{m}$. (c) Vers \ang{080.7}, il est incliné de l'autre côté~:
$\psi = \ang{30} + \ang{12.16} = \ang{42.16}$, $\cos\psi = 0.7413$, et la
carotte mesure $176.0/0.7413 = \qty{237.4}{m}$. La carotte est la plus courte,
égale à $t$, quand le sondage suit la normale, plongeant de \ang{77.84} vers
\ang{260.7}.
\end{solution}

\begin{solution}{exo:b1:cross-sections:12}
Dans le plan vertical est--ouest, avec les angles mesurés depuis l'est, les
couches supérieures montent vers l'est de \ang{20} et les couches inférieures
descendent vers l'est de \ang{70}. Faire tourner l'\omterm{def:b1:geological-maps:geological-map}{affleurement} autour de la
direction de \ang{20}, côté est vers le bas, rend les couches supérieures
horizontales et amène les couches inférieures à $\ang{70} + \ang{20} = \ang{90}$~:
elles étaient verticales quand les couches supérieures se sont déposées.
Histoire, depuis l'événement le plus ancien~: les couches inférieures se déposent
à plat~; elles sont redressées sur la tranche, jusqu'à la verticale, et
soulevées~; leurs tranches sont érodées en une surface~; les couches supérieures
se déposent en travers de cette surface~; l'ensemble est basculé de \ang{20} vers
l'ouest~; l'érosion taille la surface du sol d'aujourd'hui. C'est l'histoire de
l'\omterm{def:b1:geological-maps:geological-map}{affleurement} de Hutton à Siccar Point.
\end{solution}

\begin{solution}{pb:b1:cross-sections:1}
\textbf{1.} L'épaisseur au-dessus du tunnel est $z - 200$~: la plus grande en
$x = \qty{1}{km}$, $556 - 200 = \qty{356}{m}$~; la plus petite en $x =
\qty{9}{km}$, $339 - 200 = \qty{139}{m}$.

\textbf{2.} \qty{1}{km} d'altitude mesure \qty{4.2}{cm} et \qty{1}{km} de
distance \qty{1.4}{cm}~: $e = 4.2/1.4 = 3$.

\textbf{3.} $380.64 - 220.48 = \qty{160.16}{m}$.

\textbf{4.} En B2, la base du calcaire est à $380.64 - 100.48 =
\qty{280.16}{m}$ et la base du grès à \qty{160.16}{m}~: le tunnel, à \qty{200}{m},
est dans le grès.

\textbf{5.} B1: $524.02 - 668.98 = \qty{-144.96}{m}$; B3: $529.98 - 674.93 =
\qty{-144.95}{m}$.

\textbf{6.} D'ouest en est~: craie, argilite, calcaire, grès (au cœur de
l'\omterm{def:b1:geological-maps:fold}{anticlinal}), calcaire, argilite, et de nouveau craie jusqu'à
$x = \qty{9}{km}$, à travers le \omterm{def:b1:geological-maps:fold}{synclinal} et au-delà de la faille.

\textbf{7.} $\gamma = \ang{80.7} - \ang{90} = \ang{-9.3}$~; $\tan\beta =
\tan\ang{12.16}\cos\ang{9.3} = 0.2155\times0.9869 = 0.2126$~: $\beta =
\ang{12.0}$, vers le bas en direction de l'est.

\textbf{8.} $(280.16 + 144.95)/2000 = 425.11/2000 = 0.2126$: \ang{12.0}.

\textbf{9.} $\arctan(3\times0.2126) = \arctan 0.6377 = \ang{32.5}$.

\textbf{10.} $\gamma = \ang{120} - \ang{90} = \ang{30}$~; $\tan\beta =
\tan\ang{60}\cos\ang{30} = 1.732\times0.866 = 1.500$~: $\beta = \ang{56.3}$.

\textbf{11.} Les flancs plongent vers \ang{080.7} ou \ang{279.3}, à \ang{9.3}
seulement de la ligne de coupe, et $\cos\ang{9.3} = 0.987$~; la faille plonge
vers \ang{120}, à \ang{30} d'elle, et $\cos\ang{30} = 0.866$.

\textbf{12.} Le plan de faille est $z = 250 - \tan\ang{60}\,[(x - 4300)
\sin\ang{120} + (y - 1000)\cos\ang{120}]$. Avec $y = 6000$, $(y - 1000)
\cos\ang{120} = -2500$, donc $(x - 4300)\sin\ang{120} = 2500 - (z - 250)/
\tan\ang{60}$, et $1/\tan\ang{60} = \tan\ang{30}$.

\textbf{13.} Pour $z = 200$~: $x = 4300 + (2500 + 50\times0.5774)/0.8660 =
4300 + 2528.9/0.8660 = \qty{7220}{m}$, soit \qty{7.22}{km} de C.

\textbf{14.} \omterm{def:b1:cross-sections:throw}{Rejet horizontal} $= 200/\tan\ang{60} = \qty{115.5}{m}$~; séparation
le long de C--D $= 200/1.500 = \qty{133.3}{m}$.

\textbf{15.} La faille plonge vers l'est-sud-est et est normale~: le compartiment
supérieur, à l'est, est affaissé. Au niveau du tunnel, la craie se trouve des
deux côtés de la faille~: à l'ouest, la craie du \omterm{def:b1:geological-maps:fold}{synclinal}~; à l'est, de la craie
descendue de \qty{200}{m}.

\textbf{16.} Le tunnel reste dans la craie de part et d'autre de la faille, si
bien qu'aucun changement de roche au front ne l'annonce~; seules la signalent une
zone de roche broyée, et peut-être de l'eau. Les ingénieurs prévoient sa
position à partir de la coupe (ici \qty{7.22}{km}), forent des sondages de
reconnaissance en avant de la machine à mesure qu'elle s'en approche, et
préparent un soutènement pour la zone broyée.

\textbf{17.} $668.98 - 488.98 = \qty{180.00}{m}$.

\textbf{18.} $t = 180\cos\ang{12.16} = \qty{176.0}{m}$.

\textbf{19.} Sur la crête ($x = 3500$), la base du grès est à
\qty{160.16}{m} (question 3)~; elle descend sur les deux flancs avec le \omterm{def:b1:cross-sections:apparent-dip}{pendage
apparent} \ang{12.0} de la question 7, c'est-à-dire de $\tan\ang{12.0}$ par mètre
de distance horizontale à l'axe.

\textbf{20.} La base du calcaire est \qty{120}{m} au-dessus de la base du grès, à
$280.16 - u\tan\ang{12.0}$ avec $u = |x - 3500|$, et son toit
\qty{180}{m} plus haut, à $460.16 - u\tan\ang{12.0}$. Avec $\tan\ang{12.0} =
0.21256$~: la base est au-dessous de \qty{200}{m} lorsque
$u > 80.16/0.21256 = 377.1$, et le toit au-dessus lorsque
$u < 260.16/0.21256 = 1224.0$.

\textbf{21.} $x = 3500 \pm 377.1$ et $3500 \pm 1224.0$~: le tunnel est dans le
calcaire de \qty{2276}{m} à \qty{3123}{m} et de \qty{3877}{m} à
\qty{4724}{m}.

\textbf{22.} Le long d'une horizontale, la base et le toit du calcaire, distants
de \qty{180}{m} à la verticale et tous deux de pente $\tan\beta$, sont traversés
à une distance horizontale $h/\tan\beta$ l'un de l'autre~:
$180/0.21256 = \qty{846.8}{m}$ sur chaque flanc, soit, aux arrondis près,
$1224.0 - 377.1 = 846.9$.

\textbf{23.} $2\times846.8 = \qty{1693.7}{m}$~: le tunnel traverse
\textbf{environ \qty[text-series-to-math=true]{1694}{m} de calcaire}, plus d'un
cinquième de sa longueur de \qty{8}{km}, en deux tronçons de \qty{847}{m} de part
et d'autre du cœur de grès de l'\omterm{def:b1:geological-maps:fold}{anticlinal}.
\end{solution}
